% ECONOMICS_PROBLEM: {"id": "C73-487", "title": "Trajectory-error certificates for constrained feedback Nash approximations", "jel_code": "C73", "source_id": "OP-0229"}
% !TeX program = xelatex
\documentclass[11pt,a4paper]{article}
\usepackage[margin=23mm]{geometry}
\usepackage{fontspec}
\setmainfont{Latin Modern Roman}
\usepackage{amsmath,amssymb,mathtools}
\usepackage{xcolor,enumitem,booktabs,longtable,array,xurl,hyperref}
\definecolor{muted}{HTML}{53636C}
\hypersetup{colorlinks=true,urlcolor=blue,linkcolor=blue}
\setlength{\parindent}{0pt}
\setlength{\parskip}{5pt}
\newcommand{\R}{\mathbb R}
\newcommand{\E}{\mathbb E}
\newcommand{\N}{\mathbb N}
\newcommand{\ind}{\mathbf 1}
\newcommand{\Var}{\operatorname{Var}}
\newcommand{\Cov}{\operatorname{Cov}}
\newcommand{\supp}{\operatorname{supp}}
\DeclareMathOperator*{\argmin}{arg\,min}
\DeclareMathOperator*{\argmax}{arg\,max}
\newcommand{\entrymeta}[3]{{\sffamily\small\color{muted}\textbf{\texttt{#1}}\quad|\quad #2\par #3\par}}
\newcommand{\Problemref}[1]{\texttt{#1}}
\newcommand{\JELref}[2]{\texttt{#2} (JEL \texttt{#1})}
\begin{document}
\section*{Trajectory-error certificates for constrained feedback Nash approximations}
\textbf{Problem ID:} \texttt{C73-487}\quad \textbf{JEL:} \texttt{C73}\par
% BEGIN ECONOMICS STATEMENT
\label{problem:OP-0229}
% GAME_THEORY_PROBLEM: {"local_id": "GT-099", "original_number": 99, "kind": "R", "entry_kind": "statement_only", "published": false, "review_required": true, "category": "game_theory"}
\label{game:GT-099}
{\sffamily\small\color{muted}
\noindent\textbf{ID:} GT-099. \textbf{Category:} Differential games and model predictive control.
\textbf{Status:} R; rigorous benchmark for a source-stated extension direction.
\par}
\subsection*{Definitions and mathematical statement}
Consider an $N$-player continuous-time linear-quadratic game, with $N\geq2$,
\[
 \dot x=Ax+\sum_{i=1}^N B_i u_i,\qquad
 J_i(\mu;t_0,x_0)=\frac12\int_{t_0}^\infty
       \left(x^\top Q_i x+\sum_{j=1}^N u_j^\top R_{ij}u_j\right)dt.
\]
Here $x\in\mathbb R^d$, $u_i\in\mathbb R^{r_i}$,
$A\in\mathbb R^{d\times d}$, $B_i\in\mathbb R^{d\times r_i}$,
$Q_i\succeq0$, $R_{ij}\succeq0$, and $R_{ii}\succ0$, with symmetric cost matrices of the indicated dimensions. The nonnegativity of every $R_{ij}$ is a scope restriction of this benchmark.

Fix nonempty closed state and input constraint sets $\mathcal X\subseteq\mathbb R^d$ and $\mathcal U_i\subseteq\mathbb R^{r_i}$ and a nonempty state domain $D\subseteq\mathcal X$. Feedbacks are maps $\mu_i:[0,\infty)\times\mathcal X\to\mathcal U_i$, measurable in time and locally Lipschitz in state, for which the closed-loop equation has a unique forward solution. A profile is admissible on $D$ if, from every $(t_0,x_0)\in[0,\infty)\times D$, its solution stays in $\mathcal X$ and every $J_i$ is finite. At $(t_0,x_0)$, a unilateral deviation is feasible if, with the other feedbacks fixed, it has these properties from $(t_0,x_0)$.

A feedback Nash profile $\mu^*$ on $D$ satisfies, for every $i$, every $t_0\geq0$, every $x_0\in D$, and every feasible deviation $\nu_i$,
\[
 J_i(\mu^*;t_0,x_0)\leq J_i(\nu_i,\mu^*_{-i};t_0,x_0).
\]
Thus the Nash inequality is required at every continuation time on $D$. Require the chosen reference and candidate profiles to keep trajectories starting in $D$ within $D$, so that their reached states lie in the equilibrium domain. For an admissible approximate profile $\widehat\mu$, write
$x^{\widehat\mu}(t;x_0)$ and $x^{\mu^*}(t;x_0)$ for trajectories starting at time zero from the same initial state. The benchmark is to identify a stated class of constrained games, equilibrium selection $\mu^*$, and computable data $\eta$ (for example horizons, gain errors, stability margins, and solver certificates) for which a finite certified bound
\[
 \sup_{t\geq0}\|x^{\widehat\mu}(t;x_0)-x^{\mu^*}(t;x_0)\|
      \leq E(\eta,x_0)\qquad(x_0\in D)
\]
holds, with $E\to0$ under explicitly specified approximation limits. The assumptions must include the admissibility and stability conditions needed for this conclusion.

\subsection*{Short English statement}
Extend quantitative state-trajectory error guarantees to constrained feedback approximations, with verifiable conditions controlling active constraints and stability.

\subsection*{Variants and scope boundaries}
Constrained feedbacks may be piecewise defined or nonsmooth; a weaker regularity class must still provide a precise solution concept and uniqueness or selection. Finite-time error bounds are weaker than the displayed infinite-time bound. Equilibrium multiplicity makes the selected reference essential. A trajectory bound is neither a bound on feedback-gain distance nor a bound on profitable unilateral deviations. No residual-to-gain estimate is assumed.

\subsection*{Source relationship and status}
The source's Theorem 4 bounds the deviation of MPC and feedback-Nash \emph{state trajectories} using gain mismatch and closed-loop hypotheses. It does not prove the catalog's asserted $\operatorname{dist}(\widehat K,K^*)\leq F(\text{residual/data})$. Section 6.5 identifies extending the error analysis to constrained cases as further work. The displayed certificate problem is an editorially specified research target, not a theorem for every constrained LQ game.

\subsection*{References}
\begingroup\footnotesize\raggedright
\noindent [S1] B\'alint Varga, Karl Handwerker, Imre Rudas, and Peter Galambos, \emph{Approximate Feedback Nash Equilibria in Constrained Differential Games via Model Predictive Control with Upper Bound Guarantees} (2026), arXiv:2610.04614v1. Inspected locators: Section 3, game and feedback definitions; Section 4.4, constraints; Section 5, Theorem 4; Section 6.5, extension of the error analysis.
\url{https://arxiv.org/html/2610.04614v1}
\par\endgroup

\subsection*{Common conventions: Source mathematical conventions}

\textbf{Finite games.} For a finite set $A$, $\Delta(A)$ is its probability
simplex. Players choose independent mixed strategies in Nash equilibrium;
correlation is allowed only when explicitly part of the solution concept.
In a game with payoff functions $u_i$ and strategy profile $x$, an additive
$\varepsilon$ Nash equilibrium satisfies
\[
 u_i(x)\geq u_i(x_i',x_{-i})-\varepsilon
 \quad\text{for every player $i$ and every admissible deviation $x_i'$}.
\]
For numerical approximation, payoffs are normalized as locally specified.
An $\varepsilon$ payoff residual is distinct from distance at most
$\varepsilon$ to the set of exact equilibria. Point convergence, convergence
of empirical play, and convergence in distance to an equilibrium set are
also distinct.

\textbf{Computational representations.} Unless stated otherwise, a complexity
question takes rational data with binary encoding; polynomial time is measured
in total bit length. A fixed-accuracy polynomial algorithm is not necessarily
polynomial in $1/\varepsilon$, and neither is necessarily polynomial in
$\log(1/\varepsilon)$. Succinct games and value, demand, or sampling oracles
must declare their interfaces. Exact real solutions require an explicit
representation; an unspecified list of arbitrary real numbers is not a
finite computational output. A claimed algorithmic barrier is meaningful
only for its specified model and reductions.

\textbf{Dynamic games.} Histories, information, observation, and allowable
randomization are part of the model. A stationary strategy depends only on
the current state; a positional strategy in a deterministic graph game
chooses one action at each controlled vertex. Nash equilibrium at an initial
state and equilibrium after every history are different requirements.
An expectation of a pathwise $\liminf$ payoff, a $\liminf$ of expected averages,
a limiting expected average, and a uniform finite-horizon equilibrium are
different criteria. None is silently substituted for another.

\textbf{Allocation and mechanisms.} A complete allocation assigns every item
exactly once unless otherwise stated. Zero-valued goods, negative values,
capacity constraints, feasibility, lotteries, and copies of goods affect
fairness definitions and are specified locally. Dominant-strategy incentive
compatibility, Bayesian incentive compatibility, universal truthfulness,
and truthfulness in expectation impose different inequalities. Whenever
an objective ratio has zero denominator, its convention is stated or those
instances are excluded.

\textbf{Infinite-dimensional models.} Expectations, integrals, stochastic
equations, and optimization objectives require their local regularity,
integrability, filtrations, and admissible domains. A backward--forward
mean-field system is a boundary-value problem; it is not an initial-value
semiflow without an additional construction. Long-time, large-population,
and vanishing-noise limits require an order of limits and a convergence
topology. If a minimum need not be attained, the objective is its infimum
or supremum and the approximation requirement is stated separately.
% END ECONOMICS STATEMENT
\end{document}
